\noindent\textbf{Elección del procedimiento.} En régimen estacionario, líquido incompresible y sin fugas, conservar masa equivale a conservar caudal: $A_1V_1=A_2V_2$. Calcular primero el caudal y dividirlo entre la nueva área.

\begin{enumerate}

\item Formar el caudal en la sección conocida.

\[\fitmath{Q=A_{1}\cdot V_{1}}\]

\[\fitmath{Q=(0.05)\cdot (2)\approx 0.1\,\mathrm{m^{3}/s}}\]

\item La misma cantidad de líquido pasa por un área menor cada segundo.

\[\fitmath{V_{2}=\frac{Q}{A_{2}}}\]

\[\fitmath{V_{2}=\frac{(0.1)}{(0.01)}\approx 10\,\mathrm{m/s}}\]

\end{enumerate}
